% -------------------------------------------------------------------------
% WBS ID: RES-STR-FSI
% Deliverable: Fluid–Structure Interaction and Coupling Fidelity
% Status: In progress
% Evidence status: FSI hierarchy retained as stretch extension
% Review status: Not reviewed
% Last updated: 2026-07-19
% Dependencies: RES-STR-LOD, RES-MOD-EQS, RES-NUM-SPEC
% Downstream interfaces: Future FSI extension, RES-VER-MET, Software data schema
% -------------------------------------------------------------------------

\section{RES-STR-FSI --- Fluid–Structure Interaction and Coupling Fidelity}
\label{sec:res-str-fsi}

\subsection{Purpose and Scope}
\label{sec:res-str-fsi-purpose}

This Deliverable records a future FSI capability path and the load
history information needed to support it. The active Studentship
baseline remains fixed terrain and fixed rigid barriers in the local
URANS--VOF model. One-way CFD-to-FEA and two-way FSI are stretch
extensions after hydrodynamic load extraction is stable.

\subsection{Research Questions}
\label{sec:res-str-fsi-research-questions}

% TODO: State the questions that must be answered before the Deliverable
% can be accepted.

\subsection{Terminology and Definitions}
\label{sec:res-str-fsi-definitions}

% TODO: Define the technical terminology, symbols, quantities and
% classifications used in this Deliverable.

\subsection{Literature Evidence}
\label{sec:res-str-fsi-literature-evidence}

% TODO: Record evidence from peer-reviewed papers, authoritative datasets,
% standards, technical reports and primary documentation.
%
% Do not create a detached literature-summary list. Explain how each source
% supports, contradicts or limits a project decision.

\subsubsection{Baragamage and Wu (2024)}
\label{sec:res-str-fsi-evidence-baragamage-2024}

\textcite{baragamage2024fullycoupled} demonstrated a tsunami-impact
framework in which an incompressible free-surface Navier--Stokes model
was coupled with a structural equation of motion. Fluid pressure
generated structural force, while structural displacement and velocity
were returned to the fluid boundary.

The study establishes the feasibility of two-way tsunami FSI but does
not establish that two-way coupling is required for every barrier
candidate. Its bridge representation was structurally reduced, and
parts of the validation used a constrained flow configuration.

\subsubsection{Istrati and Buckle (2019)}
\label{sec:res-str-fsi-evidence-istrati-2019}

\textcite{istrati2019trappedair} demonstrated that the distribution of
pressure and structural demand cannot always be reconstructed from the
total force alone. Changes in internal ventilation altered uplift,
moment, and the distribution of demand between structural components.

\subsection{Governing Relations and Quantitative Definitions}
\label{sec:res-str-fsi-governing-relations}

% TODO: Record relevant equations, dimensionless groups, empirical models,
% extraction rules or quantitative definitions.
%
% Use align environments for displayed equations.
% Every displayed equation must have a unique \label{eq:...}.
% Do not place punctuation after displayed equations.

\subsubsection{Structural State and Momentum Equation}
\label{sec:res-str-fsi-structural-equations}

The structural state is provisionally defined as:

\begin{align}
  \mathcal{Q}_{\mathrm{s}}
  &=
  \left\{
    \mathbf{d}_{\mathrm{s}},
    \dot{\mathbf{d}}_{\mathrm{s}},
    \boldsymbol{\sigma}_{\mathrm{s}},
    \boldsymbol{\xi}_{\mathrm{s}}
  \right\}
  \label{eq:res-str-fsi-structural-state}
\end{align}

where $\boldsymbol{\xi}_{\mathrm{s}}$ contains the internal variables
required by the selected material model.

In continuum form, structural momentum conservation is:

\begin{align}
  \rho_{\mathrm{s}}
  \ddot{\mathbf{d}}_{\mathrm{s}}
  &=
  \nabla_{0}\cdot\mathbf{P}_{\mathrm{s}}
  +
  \rho_{\mathrm{s}}\mathbf{b}_{\mathrm{s}}
  \label{eq:res-str-fsi-structural-momentum}
\end{align}

where $\mathbf{P}_{\mathrm{s}}$ is the first
Piola--Kirchhoff stress tensor.

The corresponding finite-element semi-discrete form is:

\begin{align}
  \mathbf{M}_{\mathrm{s}}
  \ddot{\mathbf{d}}_{\mathrm{s}}
  +
  \mathbf{C}_{\mathrm{s}}
  \dot{\mathbf{d}}_{\mathrm{s}}
  +
  \mathbf{f}_{\mathrm{int}}
  \left(
    \mathbf{d}_{\mathrm{s}},
    \boldsymbol{\xi}_{\mathrm{s}}
  \right)
  &=
  \mathbf{f}_{\mathrm{fluid}}
  +
  \mathbf{f}_{\mathrm{ext}}
  \label{eq:res-str-fsi-semidiscrete-structural-equation}
\end{align}

The formulation shall support more than a single structural
degree of freedom. The initial structural model may be linear elastic,
but the interface shall permit later activation of:

\begin{itemize}
  \item geometric nonlinearity;
  \item material plasticity;
  \item contact and connection behaviour;
  \item damage and failure variables;
  \item rigid-body translation and rotation.
\end{itemize}

\subsection{Fluid--Structure Interaction Model and Activation Strategy}
\label{sec:res-str-fsi-local-coupling}

The structural formulation shall therefore preserve the spatial and
temporal hydrodynamic traction even where the optimisation initially
uses a fixed rigid barrier.

For the current project scope, this statement is a data-preservation
requirement only: fixed-barrier simulations must retain enough surface
history for later structural studies, but they do not move the barrier
or solve structural response.

\subsection{Material or Structural System}
\label{sec:res-str-fsi-material-structural-system}

% TODO: Complete this RES-STR Domain-specific subsection.

\subsection{Load Cases}
\label{sec:res-str-fsi-load-cases}

% TODO: Complete this RES-STR Domain-specific subsection.

\subsection{Constitutive Behaviour}
\label{sec:res-str-fsi-constitutive-behaviour}

% TODO: Complete this RES-STR Domain-specific subsection.

\subsection{Response and Stability Measures}
\label{sec:res-str-fsi-response-stability-measures}

% TODO: Complete this RES-STR Domain-specific subsection.

\subsection{Damage and Failure Modes}
\label{sec:res-str-fsi-damage-failure-modes}

% TODO: Complete this RES-STR Domain-specific subsection.

\subsection{Fluid--Structure Coupling Requirements}
\label{sec:res-str-fsi-fluid-structure-coupling-requirements}

% TODO: Complete this RES-STR Domain-specific subsection.

\subsubsection{Fluid--Structure Interface Conditions}
\label{sec:res-str-fsi-interface-conditions}

At the fluid--structure interface $\Gamma_{\mathrm{fs}}$, kinematic
compatibility requires:

\begin{align}
  \mathbf{u}_{\mathrm{f}}
  &=
  \dot{\mathbf{d}}_{\mathrm{s}}
  \qquad
  \text{on }
  \Gamma_{\mathrm{fs}}
  \label{eq:res-str-fsi-kinematic-compatibility}
\end{align}

Dynamic equilibrium requires:

\begin{align}
  \boldsymbol{\sigma}_{\mathrm{f}}
  \mathbf{n}_{\mathrm{f}}
  +
  \boldsymbol{\sigma}_{\mathrm{s}}
  \mathbf{n}_{\mathrm{s}}
  &=
  \mathbf{0}
  \qquad
  \text{on }
  \Gamma_{\mathrm{fs}}
  \label{eq:res-str-fsi-traction-equilibrium}
\end{align}

The structural load is obtained from the mapped fluid traction:

\begin{align}
  \mathbf{f}_{\mathrm{fluid}}
  &=
  \mathcal{T}_{\mathrm{f}\rightarrow\mathrm{s}}
  \left[
    \boldsymbol{\sigma}_{\mathrm{f}}
    \mathbf{n}_{\mathrm{f}}
  \right]
  \label{eq:res-str-fsi-fluid-to-structure-transfer}
\end{align}

where $\mathcal{T}_{\mathrm{f}\rightarrow\mathrm{s}}$ is the
fluid-to-structure load-transfer operator.

For two-way coupling, the deformed structural state is returned through:

\begin{align}
  \Gamma_{\mathrm{fs}}^{\,n+1}
  &=
  \mathcal{T}_{\mathrm{s}\rightarrow\mathrm{f}}
  \left[
    \mathbf{d}_{\mathrm{s}}^{\,n+1}
  \right]
  \label{eq:res-str-fsi-structure-to-fluid-transfer}
\end{align}

The transfer operators shall conserve resultant force and moment to
within the tolerances defined by `RES-VER-MET`.

\subsubsection{FSI Fidelity Hierarchy}
\label{sec:res-str-fsi-fidelity-hierarchy}

The structural-coupling hierarchy is:

\begin{table}[htbp]
  \centering
  \small
  \renewcommand{\arraystretch}{1.25}
  \begin{tabular}{
      p{0.11\textwidth}
      p{0.25\textwidth}
      p{0.27\textwidth}
      p{0.27\textwidth}
    }
    \hline
    Mode
    &
    Fluid treatment
    &
    Structural treatment
    &
    Intended use
    \\
    \hline
    0
    &
    Local URANS
    &
    Fixed rigid barrier
    &
    Initial fixed-rigid hydrodynamic comparison
    \\
    1
    &
    Local URANS or LES
    &
    One-way CFD-to-FEA
    &
    Stress, deformation, stability, and failure screening
    \\
    2
    &
    Local URANS or LES
    &
    Partitioned two-way FSI
    &
    Cases with material but moderate deformation feedback
    \\
    3
    &
    Local high-fidelity model
    &
    Strongly coupled nonlinear FSI
    &
    Flexible, moving, unstable, or intentionally compliant final
    designs
    \\
    \hline
  \end{tabular}
  \caption{Fluid--structure coupling hierarchy for candidate and
  shortlisted barrier designs}
  \label{tab:res-str-fsi-fidelity-hierarchy}
\end{table}

\subsubsection{Two-Way FSI Activation Gate}
\label{sec:res-str-fsi-activation-gate}

Two-way FSI shall be activated where one or more of the following
conditions are satisfied:

\begin{enumerate}
  \item structural displacement changes the exposed geometry or
    wetted area materially;
  \item structural motion changes a flow passage, opening, or
    overtopping path;
  \item structural velocity is no longer negligible relative to the
    adjacent fluid velocity;
  \item rotation or translation changes the reflected or transmitted
    wave field;
  \item connection flexibility controls the global response;
  \item instability, uplift, sliding, overturning, or large rotation
    is predicted;
  \item the barrier relies on intentional compliant motion for energy
    dissipation;
  \item one-way and two-way pilot simulations produce materially
    different hydrodynamic or structural performance measures.
\end{enumerate}

A preliminary motion-feedback indicator may be defined as:

\begin{align}
  \Pi_{\mathrm{fsi}}
  &=
  \max_{\Gamma_{\mathrm{fs}},t}
  \left(
    \frac{
      \left\lVert
      \dot{\mathbf{d}}_{\mathrm{s}}
      \right\rVert
    }{
      \left\lVert
      \mathbf{u}_{\mathrm{f}}
      \right\rVert
      +
      \epsilon_{\mathrm{u}}
    }
  \right)
  \label{eq:res-str-fsi-motion-feedback-indicator}
\end{align}

where $\epsilon_{\mathrm{u}}$ prevents division by zero.

The threshold for activating two-way FSI shall be determined through
sensitivity testing rather than assigned universally.

\subsection{Reference Descriptions of the Coupled Domains}
\label{subsec:res_str_fsi_reference_descriptions}

The regional and local fluid models use an Eulerian description in which fluid variables are evaluated at spatial positions

\begin{align}
\bm{u}
&=
\bm{u}\!\left(\bm{x},t\right)
\label{eq:res_str_fsi_eulerian_velocity}
\end{align}

The solid uses a Lagrangian material description

\begin{align}
\bm{x}
&=
\bm{\chi}\!\left(\bm{X},t\right)
=
\bm{X}
+
\bm{d}\!\left(\bm{X},t\right)
\label{eq:res_str_fsi_lagrangian_motion}
\end{align}

Impact does not require a Lagrangian fluid description. Retaining an Eulerian fluid mesh is more suitable for breaking, overtopping and free-surface topology change, while the Lagrangian solid description tracks the barrier deformation and material history.

\subsection{Fluid--Solid Interface Conditions}
\label{subsec:res_str_fsi_interface_conditions}

Two-way coupling requires kinematic compatibility and dynamic equilibrium on the moving interface

\begin{align}
\bm{u}_{f}
&=
\dot{\bm{d}}_{s}
\qquad
\text{on }\Gamma_{fs}(t)
\label{eq:res_str_fsi_kinematic_condition}
\\
\bm{\sigma}_{f}\bm{n}_{f}
+
\bm{\sigma}_{s}\bm{n}_{s}
&=
\bm{0}
\qquad
\text{on }\Gamma_{fs}(t)
\label{eq:res_str_fsi_dynamic_condition}
\end{align}

The fluid solver transfers traction to the structure. The structural solver returns the current interface position and velocity to the fluid representation. A one-way calculation omits the return of structural motion, while a two-way calculation repeats this exchange throughout the transient solution.

\subsection{Moving-Boundary Treatment}
\label{subsec:res_str_fsi_moving_boundary_treatment}

A body-fitted Arbitrary Lagrangian--Eulerian method moves the fluid mesh with the structure. The fluid advection velocity is then evaluated relative to the mesh velocity

\begin{align}
\bm{u}_{\mathrm{rel}}
&=
\bm{u}
-
\bm{w}_{m}
\label{eq:res_str_fsi_ale_relative_velocity}
\end{align}

ALE provides a direct conforming interface and conventional near-wall treatment. Its principal limitation for the present application is the risk of fluid-mesh distortion under severe structural deformation, followed by remeshing and transfer of pressure, velocity, VOF and turbulence fields. It remains useful as a reference method for moderate-deformation benchmarks \autocite{girfoglio2021}

If FSI is activated as a stretch extension, the selected reference
family is a sharp immersed structural surface with conservative cut
cells on a fixed Eulerian background mesh. The method avoids global
fluid-mesh deformation while accounting geometrically for the changing
fluid volume and moving solid boundary. Pasquariello et al.
demonstrate conservative finite-volume--finite-element coupling with
cut cells and nonmatching interface transfer for three-dimensional FSI
with large structural deformation \autocite{pasquariello2016}.
Baragamage and Wu apply a related immersed-boundary and cut-cell
strategy to three-dimensional tsunami loading
\autocite{baragamage2024fullycoupled}

\textbf{Deferred selected approach.} Retain a sharp immersed-boundary
representation with conservative cut-cell treatment as the preferred
FSI extension path. Retain ALE as a comparator for
moderate-deformation verification cases.

The selected family must subsequently define cut-cell geometry, changing fluid volumes, small-cell stabilisation, reconstructed wall states, traction recovery and conservative transfer between fluid and structural interfaces.

\subsection{Deferred High-Level FSI Architecture}
\label{subsec:res_str_fsi_current_architecture}

The deferred high-level architecture is

\begin{align}
&\text{Eulerian polyhedral finite-volume two-phase fluid}
\nonumber\\
&\quad+
\text{VOF free-surface capture}
\nonumber\\
&\quad+
\text{Updated-Lagrangian three-dimensional solid FEM}
\nonumber\\
&\quad+
\text{generalised-}\alpha\text{ structural time integration}
\nonumber\\
&\quad+
\text{Newton-type nonlinear structural solution}
\nonumber\\
&\quad+
\text{sharp immersed boundary with conservative cut cells}
\label{eq:res_str_fsi_current_architecture}
\end{align}

The following high-level decisions remain unresolved:

\begin{enumerate}
    \item monolithic or partitioned coupling
    \item loose or strong partitioned coupling if the partitioned family is selected
    \item added-mass stabilisation, relaxation and interface quasi-Newton acceleration
    \item conservative transfer of force, moment, displacement and interface work
    \item coordination of fluid and structural time levels
    \item transition criteria between rigid, one-way and two-way FSI calculations
\end{enumerate}

\subsection{Candidate Approaches}
\label{sec:res-str-fsi-candidate-approaches}

% TODO: Define the credible candidate approaches considered.

\subsection{Critical Comparison}
\label{sec:res-str-fsi-critical-comparison}

% TODO: Compare candidates using physically and project-relevant criteria.
% Do not select an approach solely on popularity or implementation ease.

\subsection{Assumptions and Applicability}
\label{sec:res-str-fsi-assumptions}

% TODO: State assumptions and the conditions under which they are valid.

\subsection{Limitations and Failure Conditions}
\label{sec:res-str-fsi-limitations}

% TODO: State where the evidence, model or selected approach ceases to be
% reliable or applicable.

\subsection{Project-Specific Conclusions}
\label{sec:res-str-fsi-conclusions}

% TODO: Convert the evidence into conclusions specific to the Studentship
% project. Do not merely restate literature findings.

\subsubsection{Project-Specific FSI Decision}
\label{sec:res-str-fsi-project-decision}

The software and data architecture shall be FSI-capable from the
outset. The baseline hydrodynamic comparison model shall nevertheless use a fixed
rigid barrier and retain complete surface traction histories.

The initial structural verification shall use one-way CFD-to-FEA load
transfer. Two-way FSI shall be introduced only where the activation
gate demonstrates that structural motion feeds back materially into
the hydrodynamic response.

This approach retains the equation-level capability for coupled
simulation without imposing fully coupled FSI on the active barrier
comparison campaign.

\subsection{Selected Approach and Rationale}
\label{sec:res-str-fsi-selected-approach}

% TODO: Record the selected approach only after sufficient evidence exists.
% State the alternatives rejected and the reasons for rejection.

\subsection{Unresolved Decisions}
\label{sec:res-str-fsi-unresolved-decisions}

% TODO: Record unresolved questions, required evidence and decision owners.

\subsection{Requirements and Handoffs}
\label{sec:res-str-fsi-requirements}

Software shall preserve load, geometry and patch histories in a form
that a later FSI extension can consume, but shall not implement moving
boundaries, structural time integration or two-way coupling from this
Research update. RES-VER-MET shall define force/moment transfer
consistency metrics before FSI activation.

\subsection{Deliverable Completion Record}
\label{sec:res-str-fsi-completion-record}

% TODO: Record whether the Deliverable is:
%
% - Not started
% - In progress
% - Draft complete
% - Reviewed
% - Accepted
%
% Acceptance requires:
%
% - the research questions to be answered;
% - claims to be supported by appropriate evidence;
% - assumptions and limitations to be explicit;
% - the project-specific conclusion to be stated;
% - downstream requirements to be traceable;
% - unresolved decisions to be closed or formally deferred.
